% UTAC v2.0 - ENDORSEMENT REQUEST LETTER TEMPLATE
% Professional template for contacting potential endorsers

\documentclass[11pt,a4paper]{article}

\usepackage[utf8]{inputenc}
\usepackage[margin=2.5cm]{geometry}
\usepackage{hyperref}

% No page numbers
\pagestyle{empty}

% ============================================================
% DOCUMENT
% ============================================================
\begin{document}

% ============================================================
% LETTERHEAD
% ============================================================
\begin{flushleft}
\textbf{Johann Römer}\\
Independent Researcher\\
Marburg, Germany\\
\texttt{[Contact via Zenodo: 10.5281/zenodo.17472834]}\\
\end{flushleft}

\vspace{1em}

\begin{flushleft}
\today
\end{flushleft}

\vspace{1em}

\begin{flushleft}
\textbf{[RECIPIENT NAME]}\\
[RECIPIENT TITLE]\\
[RECIPIENT INSTITUTION]\\
[RECIPIENT EMAIL]
\end{flushleft}

\vspace{1em}

% ============================================================
% SUBJECT LINE
% ============================================================
\textbf{Subject: arXiv Endorsement Request / Collaboration Inquiry — UTAC Framework (Threshold Dynamics)}

\vspace{1em}

% ============================================================
% LETTER BODY
% ============================================================
Dear Professor [LAST NAME],

\vspace{0.5em}

% PARAGRAPH 1: Introduction + Why Them
I am writing to request your endorsement for submitting my manuscript \textit{"Domain-Specific Universality in Threshold Activation Criticality: A Multi-Attractor Framework"} to arXiv (\texttt{physics.data-an} or \texttt{cond-mat.stat-mech}). 

[PERSONALIZATION BLOCK — Choose relevant option:]

\textbf{[OPTION A — Their Research Aligns:]}\\
Your research on [SPECIFIC TOPIC, e.g., "early warning signals for climate tipping points" / "emergent abilities in large language models" / "critical phenomena in complex systems"] directly intersects with my findings. Specifically, your [PAPER TITLE, YEAR] demonstrated [KEY RESULT], which aligns remarkably with the UTAC framework's prediction that [CONNECTION TO THEIR WORK].

\textbf{[OPTION B — Methodological Connection:]}\\
Your expertise in [RENORMALIZATION GROUP THEORY / STATISTICAL MECHANICS / COMPLEX SYSTEMS MODELING] makes you ideally positioned to evaluate this work. The UTAC framework extends Wilson-Kogut RG theory to empirical threshold systems, revealing domain-specific universality classes rather than strict convergence.

\textbf{[OPTION C — Shared Interest in Applications:]}\\
Given your work on [AI SAFETY / CLIMATE TIPPING POINTS / CONSCIOUSNESS / NEURODEGENERATIVE DISEASES], you may find particular interest in the UTAC framework's applications to [RELEVANT DOMAIN]. The discovery that [KEY FINDING RELEVANT TO THEM] has direct implications for [THEIR RESEARCH AREA].

\vspace{0.5em}

% PARAGRAPH 2: The Work (Brief Summary)
The manuscript presents empirical analysis of 78 threshold systems across 5 scientific domains (Informational, Biological, Climate, Neurodegenerative, Geophysical), quantifying transition steepness via parameter $\beta$. Key findings include:

\begin{itemize}
\item \textbf{Statistical validation} of domain-specific $\beta$-clustering (ANOVA: $F(4,73) = 185.3$, $p < 10^{-20}$, $\eta^2 = 0.91$)
\item Identification of the \textbf{Computational Criticality Universality Class} (CCUC) at RG fixed point $\beta \approx 4.2$
\item Discovery of \textbf{Golden Ratio hierarchical scaling}: $\beta_n \approx \Phi^{n/3}$ validated with $< 8\%$ error
\item Applications to AI emergence (LLMs), consciousness (neuronal avalanches), and climate tipping points
\end{itemize}

The full manuscript (25 pages), dataset (78 systems, CSV format), and analysis code are publicly available on Zenodo (DOI: 10.5281/zenodo.17472834).

\vspace{0.5em}

% PARAGRAPH 3: The Request
\textbf{Request:} I would be honored if you could provide arXiv endorsement for \texttt{physics.data-an} (primary) or \texttt{cond-mat.stat-mech} (secondary). I am also open to discussing potential collaboration or co-authorship if the work aligns with your research interests.

[OPTIONAL — If seeking journal submission advice:]\\
Alternatively, if you believe this work would benefit from further development before arXiv submission, I would greatly value your guidance on next steps. I am targeting \textit{Chaos} or \textit{Physical Review E} for peer review.

\vspace{0.5em}

% PARAGRAPH 4: Closing
Thank you for considering this request. I have attached the executive summary (2 pages) for quick reference, with the full manuscript available at the Zenodo link above. I am happy to provide any additional information or clarification you may need.

\vspace{1em}

Sincerely,

\vspace{2em}

Johann Römer\\
Independent Researcher\\
\texttt{Zenodo: https://doi.org/10.5281/zenodo.17472834}

% ============================================================
% ATTACHMENTS
% ============================================================
\vspace{2em}

\textbf{Attachments:}
\begin{itemize}
\item Executive Summary (2 pages)
\item Link to Full Manuscript \& Dataset: \url{https://doi.org/10.5281/zenodo.17472834}
\end{itemize}

% ============================================================
% FOOTER WITH KEY LINKS
% ============================================================
\vspace{2em}
\hrule
\vspace{0.5em}

\begin{center}
\small
\textbf{Quick Access:}\\
Zenodo Repository: \url{https://doi.org/10.5281/zenodo.17472834}\\
Main Paper (PDF): [Direct link after Zenodo upload]\\
Executive Summary (PDF): [Direct link after Zenodo upload]
\end{center}

\end{document}
